\documentclass[10pt,openany]{book}
\usepackage[utf8]{inputenc}
\usepackage[a4paper,hmargin={4cm,4cm},vmargin={3.5cm,3.5cm}]{geometry}
\usepackage[colorlinks=true,bookmarksopen=false,linkcolor=blue,citecolor=red]{hyperref}
\usepackage{subfiles}
\usepackage{calc} 
\usepackage{datetime}
\usepackage{comment}
\usepackage{amssymb}
\usepackage{amsthm}
\usepackage{amsmath}
\usepackage{amsrefs}
\usepackage{titlesec}
\usepackage{titletoc}
\usepackage{dsfont}
\usepackage{euscript}
\usepackage{fourier-orns}
%\usepackage{pxfonts}
%\usepackage{newpxtext}
%\usepackage{tgpagella}
\usepackage{palatino}
\usepackage[sc]{mathpazo} % add possibly `sc` and `osf` options
%\usepackage{eulervm}
%\usepackage[adobe-utopia]{mathdesign}
\usepackage[T1]{fontenc}
\usepackage{graphicx}
% \usepackage{tikz}
% \usetikzlibrary{tikzmark}
\usepackage{tikz-cd}
\tikzcdset{
arrow style=tikz,
diagrams={>=latex}
}


\linespread{1.1}
\setlength{\parindent}{0ex}
\setlength{\parskip}{.4\baselineskip}
\definecolor{blue}{rgb}{0, 0.1, 0.6}

\DeclareFontFamily{OT1}{pzc}{}
\DeclareFontShape{OT1}{pzc}{m}{it}{<-> s * [1.10] pzcmi7t}{}
\DeclareMathAlphabet{\mathpzc}{OT1}{pzc}{m}{it}


\newcommand{\mylabel}[1]{{\ssf{#1}}\hfill}
\renewenvironment{itemize}
  {\begin{list}{$\triangleright$}{%
   \setlength{\parskip}{0mm}
   \setlength{\topsep}{.4\baselineskip}
   \setlength{\rightmargin}{0mm}
   \setlength{\listparindent}{0mm}
   \setlength{\itemindent}{0mm}
   \setlength{\labelwidth}{3ex}
   \setlength{\itemsep}{.4\baselineskip}
   \setlength{\parsep}{0mm}
   \setlength{\partopsep}{0mm}
   \setlength{\labelsep}{1ex}
   \setlength{\leftmargin}{\labelwidth+\labelsep}
   \let\makelabel\mylabel}}{%
   \end{list}\vspace*{-\parskip}}

%\def\emph#1{\textcolor{blue}{\boldmath\textbf{#1}}}

% \setlength{\fboxsep}{0.3mm}
% \def\emph#1{\fcolorbox{white}{gray!50}{\vphantom{Iq}#1}}
\def\emph#1{\textcolor{blue}{\textbf{\boldmath #1}}}

\def\jj{\mbox{-}}
\def\E{\exists}
\def\A{\forall}
\def\mdot{\mathord\cdot}
\def\models{\vDash}
\def\pmodels{\mathrel{\models\kern-1.5ex\raisebox{.5ex}{*}}}
\def\notmodels{\nvDash}
\def\proves{\vdash}
\def\notproves{\nvdash}
\def\proves{\vdash}
\def\provesT{\mathrel{\mathord{\vdash}\hskip-1.13ex\raisebox{-.5ex}{\tiny$T$}}}
\def\ZZ{\mathds Z}
\def\NN{\mathds N}
\def\QQ{\mathds Q}
\def\RR{\mathds R}
\def\BB{\mathds B}
\def\CC{\mathds C}
\def\DD{\mathds D}
\def\PP{\mathds P}
\def\Ar{{\rm Ar}}
\def\dom{\mathop{\rm dom}}
\def\range{\mathop{\rm img}}
\def\rank{\mathop{\rm rank}}
\def\dcl{\mathop{\rm dcl}}
\def\acl{\mathop{\rm acl}}
\def\fp{{\rm fp}}
\def\cf{\mathop{\rm cf}}
\def\rad{\mathop{\rm rad}}
\def\eq{{{\rm eq}}}
\def\ccl{{\rm ccl}}
\def\Th{\textrm{Th}}
\def\Diag{{\rm Diag}}
\def\atpmTh{\textrm{Th}_\atpm}
\def\Mod{\mathop{\rm Mod}}
\def\Rmod{{\mbox{\scriptsize $R$-mod}}}
\def\Aut{\textrm{Aut\kern.15ex}}
\def\Autf{\mathord{\rm Aut\kern.15ex{f}\kern.15ex}}
\def\Cb{\textrm{Cb\kern.15ex}}
\def\orbit{\O}
\def\oorbit{\mathpzc{o}}
\def\oorbitf{\mathpzc{of\!}}
\def\id{\textrm{id}}
\def\tp{{\rm tp}}
\def\atpm{{\tiny\rm at^\pm}}
\def\qftp{\textrm{qf\mbox{-}tp}}
\def\attp{\textrm{at\mbox{-}tp}}
\def\atpmtp{\atpm\mbox{-}\textrm{tp}}
\def\Deltatp{\Delta\mbox{-}\textrm{tp}}
\def\pmDelta{\Delta\hskip-.3ex\raisebox{1ex}{\tiny$\pm$}}
\def\pmDeltatp{\noindent\pmDelta\hskip-.3ex\textrm{-tp}}
\def\EMtp{\textrm{{\small EM}\mbox{-}tp}}

\def\nonfork{\mathop{\raise0.2ex\hbox{\ooalign{\hidewidth$\vert$\hidewidth\cr\raise-0.9ex\hbox{$\smile$}}}}}

%\def\nonforkc{\mathbin{\raise1.2ex\rlap{\kern1.3ex$\cdot$}\rlap{\kern1.1ex\rule{0.1ex}{1.9ex}}\raise-0.3ex\hbox{$\smile$} } }
%\def\cnonfork{\mathbin{\raise1.3ex\rlap{\kern0.4ex$\cdot$}\raise0.6ex\rlap{\kern0.85ex$\vert$}\raise-0.3ex\hbox{$\smile$}} }

\def\cnonfork{\mathbin{\raise1.8ex\rlap{\kern0.6ex\rule{0.6ex}{0.1ex}}\rlap{\kern1.1ex\rule{0.1ex}{1.9ex}}\raise-0.3ex\hbox{$\smile$} } }

\def\nonforkc{\mathbin{\raise1.8ex\rlap{\kern1.1ex\rule{0.6ex}{0.1ex}}\rlap{\kern1.1ex\rule{0.1ex}{1.9ex}}\raise-0.3ex\hbox{$\smile$} } }


\def\cpaw{\mathbin{\ooalign{\kern-0.4ex$-$\hidewidth\cr$<$}}}
\def\cpawdot{\ooalign{$\kern1.2ex\cdot$\cr$\cpaw$\cr}}
\def\cev#1{\reflectbox{\ensuremath{\vec{\reflectbox{\ensuremath{#1}}}}}}

\def\sm{\smallsetminus}
\def\atpmL{L_{\atpm}\hskip-.3ex}
\def\qfL{L_{\rm qf}}
\def\atL{L_{\tiny\rm at}\hskip-.3ex}
\def\simdiff{\triangle}
\def\IMP{\Rightarrow}
\def\PMI{\Leftarrow}
\def\IFF{\Leftrightarrow}
\def\NIFF{\nLeftrightarrow}
\def\imp{\rightarrow}
\def\pmi{\leftarrow}
\def\iff{\leftrightarrow}
\def\niff{\mathrel{{\leftrightarrow}\llap{\raisebox{-.1ex}{{\small$/$}}\hskip.5ex}}}
\def\nequiv{\mathrel{\mbox{$\equiv$\llap{{\small$/$}\hskip.3ex}}}}
\def\equivEM{\stackrel{\smash{\scalebox{.5}{\rm EM}}}{\equiv}}
\def\equivL{\stackrel{\smash{\scalebox{.5}{\rm L}}}{\equiv}}
\def\equivKP{\stackrel{\smash{\scalebox{.5}{\rm KP}}}{\equiv}}
\def\equivSh{\stackrel{\smash{\scalebox{.5}{\rm Sh}}}{\equiv}}
\def\dIFF{\IFF\hskip-2.2ex\smash{\raisebox{1.3ex}{\tiny def}}}
\def\deq{\mathrel{=\hskip-1.9ex\smash{\raisebox{1.2ex}{\tiny def}}}}
\def\swedge{\mathbin{\raisebox{.2ex}{\tiny$\mathbin\wedge$}}}
\def\svee{\mathbin{\raisebox{.2ex}{\tiny$\mathbin\vee$}}}

\def\bigsum{\mathop{\mbox{\large$\displaystyle\sum$}}}
\def\bigint{\mathop{\mbox{\large$\displaystyle\int$}}}


%\def\hookdoubleheadrightarrow{\hookrightarrow\mathrel{\mspace{-15mu}}\rightarrow}

\def\isomap{\rlap{\kern0.8ex\raisebox{.8ex}{\scriptsize$\kern.2ex\sim$}}\rightarrow}

\def\P{\EuScript P}
\def\D{\EuScript D}
\def\Aa{\EuScript A}
\def\Ee{\EuScript E}
\def\X{\EuScript X}
\def\Y{\EuScript Y}
\def\Z{\EuScript Z}
\def\C{\EuScript C}
\def\U{\EuScript U}
\def\Hh{\EuScript H}
\def\I{\EuScript I}
\def\V{\EuScript V}
\def\W{\EuScript W}
\def\R{\EuScript R}
\def\F{\EuScript F}
\def\G{\EuScript G}
\def\B{\EuScript B}
\def\M{\EuScript M}
\def\Ll{\EuScript L}
\def\K{\EuScript K}
\def\O{\EuScript O}
\def\J{\EuScript J}
\def\S{\EuScript S}
\def\<{\langle}
\def\>{\rangle}
\def\0{\varnothing}
\def\theta{\vartheta}
\def\phi{\varphi}
\def\epsilon{\varepsilon}
\def\ssf#1{\textsf{\small #1}}



\titlecontents{section}
[3.8em] % ie, 1.5em (chapter) + 2.3em
{\vskip-1ex}
{\contentslabel{1.5em}}
{\hspace*{-2.3em}}
{\titlerule*[1pc]{}\contentspage}


\titleformat
{\chapter} % command
[display] % shape
{\bfseries\LARGE} % format
{Chapter \ \thechapter} % label
{0.5ex} % sep
{} % before-code
[] % after-code
\titleformat{\section}[block]{\Large\bfseries}{\makebox[5ex][r]{\textbf{\thesection}}}{1.5ex}{}
\titlespacing*{\chapter}{0em}{.5ex plus .2ex minus .2ex}{2.3ex plus .2ex}
\titlespacing*{\section}{-9.7ex}{3ex plus .5ex minus .5ex}{1ex plus .2ex minus .2ex}

\renewcommand*\thesection{\arabic{section}}

\newtheoremstyle{mio}% name
     {2\parskip}%      Space above
     {\parskip}%      Space below
     {\sl}%         Body font
     {}%         Indent amount (empty = no indent, \parindent = para indent)
     {\bfseries}% Thm head font
     {}%        Punctuation after thm head
     {1ex}%     Space after thm head: " " = normal interword space;
           %   \newline = linebreak
     {\llap{\thmnumber{#2}\hskip2mm}\thmname{#1}\thmnote{\bfseries{} #3}}% Thm head spec (can be left empty, meaning `normal')

\newtheoremstyle{liscio}% name
     {2\parskip}%      Space above
     {0mm}%      Space below
     {}%         Body font
     {}%         Indent amount (empty = no indent, \parindent = para indent)
     {\bfseries}% Thm head font
     {}%        Punctuation after thm head
     {1.5ex}%     Space after thm head: " " = normal interword space;
           %   \newline = linebreak
     {\llap{\thmnumber{#2}\hskip2mm}\thmname{#1}\thmnote{\bfseries{} #3}}% Thm head spec (can be left empty, meaning `normal')



\newcounter{thm}[chapter]

\renewcommand{\thethm}{\thechapter.\arabic{thm}}

\theoremstyle{mio}
\newtheorem{theorem}[thm]{Theorem}
\newtheorem{corollary}[thm]{Corollary}
\newtheorem{proposition}[thm]{Proposition}
\newtheorem{lemma}[thm]{Lemma}
\newtheorem{fact}[thm]{Fact}
\newtheorem{definition}[thm]{Definition}
\newtheorem{assumption}[thm]{Assumption}
\newtheorem{void_thm}[thm]{}
\theoremstyle{liscio}
\newtheorem{void_def}[thm]{}
\newtheorem{remark}[thm]{Remark}
\newtheorem{notation}[thm]{Notation}
\newtheorem{note}[thm]{Note}
\newtheorem{exercise}[thm]{Exercise}
\newtheorem{example}[thm]{Example}
\setlength{\partopsep}{0mm}
\setlength{\topsep}{0mm}

\def\QED{\noindent\nolinebreak[4]\hfill\rlap{\ \ $\Box$}\medskip}
\renewenvironment{proof}[1][Proof]%
{\begin{trivlist}\item[\hskip\labelsep {\bf #1}]}
{\QED\end{trivlist}}

\newenvironment{void}[1][]%
{\begin{trivlist}\item[\hskip\labelsep {\bf #1}]}
{\QED\end{trivlist}}



\pagestyle{plain}

\definecolor{violet}{RGB}{115, 0, 205}
\definecolor{brown}{RGB}{150, 50, 10}
\definecolor{green}{RGB}{5,110, 35}
\definecolor{emphcolor}{rgb}{.98,.98,.70}

\def\mr{\color{brown}}
\def\gr{\color{green}}
\def\vl{\color{violet}}

\def\mrA{{\mr\Aa}}
\def\mrB{{\mr\B}}
\def\mrC{{\mr\C}}
\def\mrD{{\mr\D}}
\def\mrG{{\mr\G}}
\def\mrU{{\mr\U}}
\def\mrV{{\mr\V}}
\def\mrW{{\mr\W}}
\def\grB{{\gr\B}}
\def\grC{{\gr\C}}
\def\grD{{\gr\D}}
\def\grV{{\gr\V}}
\def\grCC{{\gr\CC}}
\def\grDD{{\gr\DD}}


\renewcommand*{\emph}[1]{%
\smash{\tikz[baseline]\node[rectangle, fill=emphcolor, rounded corners, inner xsep=0.2ex, inner ysep=0.2ex, anchor=base, minimum height = 2.7ex]{#1};}}


\begin{document}

\def\medrel#1{\parbox[t]{6ex}{\hfil$#1$}}
\def\ceq#1#2#3{\parbox{25ex}{$\displaystyle #1$}\medrel{#2}$\displaystyle  #3$}


Let $I\subseteq |x|=|z|$ and let $y$ be a tuple of length $|I|$. Define

\ceq{\hfill\theta_I(y,x,z)}{=}{y=x_{\restriction I}\wedge\bigwedge_{i\notin I} x_i\neq z_i}

$\bar x=\<x_i:i<\omega\>$ sequence of indiscernibles over $\Aa$.

$\theta_I(y,x_0,x_1)$

$\Aa=\phi(\U,x_0)$
\end{document}


\section{Approximations}
\def\uuu{{\mathds 1}}
\def\gruuu{{\mathds 1}}

For  $p\in S_{\mr x}(\U)$, a global type, we write ${\gr\uuu_{\phi,p}}:\U^{|{\gr z}|}\to\RR$ for the indicator function of the set ${\gr\D_{p,\phi}}$.
%
For $n\le\omega$ let $\bar p=\<p_i:i<n\>$ be a tuple of global types.
%
We write  ${\gr\uuu_{\phi,\bar p}}$ for the corresponding tuple of indicator functions.
%
For $\bar\alpha\in\RR^{n}$ we write

\ceq{\hfill\bar\alpha\cdot{\gr\uuu_{\phi,\bar p}}}{=}{\sum^n_{i=0}\alpha_i\cdot{\gr\uuu_{\phi, p_i}}}

Finally, if ${\mr b}\in\U^{|{\mr x}|}$ then ${\gr\uuu_{\phi, b}}$ denotes the indicator functon of the set $\phi({\mr b};\,{\gr\U})$.
%
Note that this coincides with ${\gr\uuu_{\phi, p}}$, for $p$ the unique global type realized in ${\mr b}$.

We say that the function $\grDD:\U^{|{\gr z}|}\to\RR$ is \emph{externally definable\/} if $\grDD=_\epsilon\bar\alpha\cdot{\gr\uuu_{\phi,\bar p}}$ for some $\bar\alpha$ and $\bar p$.

\begin{definition}\label{def_epprox}
We say that $\grDD$ is \emph{approximated\/} by the formula $\phi({\mr x}\,;{\gr z})$ if for every positive $\epsilon$ there is a finite tuple ${\mr\bar b}\in\big(\U^{|{\mr x}|}\big)^{<\omega}$ such that 

\ceq{\hfill\big\|\grDD\ -\ \bar\alpha\cdot{\gr\uuu_{\phi,\bar b}}\big\|}{<}{\epsilon,} 

where the norm is the sup-norm.
%
We may also write \emph{$\grDD=_\epsilon\bar\alpha\cdot{\gr\uuu_{\phi,\bar p}}$.}
%
If in addition we have that $\bar\alpha\cdot{\gr\uuu_{\phi,\bar b}}\le\grDD$, we say that  $\grDD$ is \emph{approximated from below}.
Similarly, when $\grDD\leq\bar\alpha\cdot{\gr\uuu_{\phi,\bar b}}$ we say that  $\grDD$ is \emph{approximated from above}.
We may call $\phi({\mr x}\,;{\gr z})$ the \emph{sort} of $\grDD$.\QED
\end{definition} 
 
The following proposition is clear by compactness.

\begin{proposition}\label{lem_approx=external}
For every $\grD$ the following are equivalent
\begin{itemize}
\item[1.] $\grD$ is approximated by $\phi({\mr x}\,;{\gr z})$;
\item[2.] $\grD$ is externally definable by $\phi({\mr x}\,;{\gr z})$.\QED
\end{itemize}
\end{proposition}

Approximability from below is an adaptation to our context of the notion of \textit{having an honest definition} in \cite{CS}.
\begin{definition}\label{def_defble_type}
We say that the global type $p\in S_{{\mr x}}(\U)$ is \emph{honestly definable\/} if for every $\phi({\mr x}\,;{\gr z})\in L$ the set ${\gr\D_{p,\phi}}$ is approximated from below (by some formula).
We say that $p$ is \emph{definable\/} if the sets ${\gr\D_{p,\phi}}$ are all definable (over $\U$).
Note that the terminology is misleading: honestly definable is weaker than definable.
\end{definition}

\begin{example}
Every definable set is trivially approximable.
Sets may be approximable by different formulas.
For instance, if $T=T_{\rm dlo}$, then $\D=\{z\in\U:a\le z\le b\}$ is approximable both from below and from above by the formula $x_1<z<x_2$ though it is not definable by this formula.

Now, let $T=T_{\rm rg}$.
Then every $\D\subseteq\U$ is approximable and, when $\D$ has small infinite cardinality, it is approximable from above but not from below, see Exercise~\ref{ex_rg_small_def_set}.\QED
\end{example}

In Definition~\ref{def_epprox}, the sort $\phi({\mr x}\,;{\gr z})$ is fixed (otherwise any set would be approximable) but this requirement of uniformity may be dropped if we allow $B$ to have larger cardinality.

\begin{proposition}\label{lem_approx_nonunif}
For every $\grD$ the following are equivalent
\begin{itemize}
\item[1.] $\grD$ is approximable;
\item[2.] for every $C\subseteq\U$ of cardinality $\le|T|$ there is $\psi({\gr z})\in L(\U)$ such that $\psi({\gr\U})=_C\grD$.
\end{itemize}
Similarly, the following are equivalent
\begin{itemize}
\item[3.] $\grD$ is approximable from below;
\item[4.]  for every ${\gr C}\subseteq\grD$ of cardinality $\le|T|$ there is $\psi({\gr z})\in L(\U)$ such that ${\gr C}\subseteq \psi({\gr\U})\subseteq\grD$.
\end{itemize}
\end{proposition}

\begin{proof}
To prove \ssf{2}$\IMP$\ssf{1} assume \ssf{2} and negate \ssf{1} for a contradiction.
For each formula $\psi({\mr x}\,;{\gr z})\in L$ choose a finite set $B$ such that $\psi({\mr b}\,;{\gr\U})\neq_B\grD$ for every ${\mr b}\in\U^{|{\mr x}|}$.
Let $C$ be the union of all these finite sets.
Clearly $|C|\le|T|$.
By \ssf{2} there are a formula $\phi({\mr x}\,;{\gr z})$ and a tuple ${\mr c}$ such that $\phi({\mr c}\,;{\gr\U})=_C\grD$, contradicting the definition of $C$.

Implication \ssf{1}$\IMP$\ssf{2} is obtained by compactness and the equivalence \ssf{3}$\IFF$\ssf{4} is proved similarly.
\end{proof}



\begin{remark}\label{prop_approx_el_eq}
If $\grD\subseteq\U^{|{\gr z}|}$ is approximated by $\phi({\mr x}\,;{\gr z})$ then so is any $\grC$ such that $\grC\equiv\grD$, see Section~\hyperref[expansions]{\ref*{invariantL}.\ref*{expansions}} for the notation.
In fact, if the set $\grD$ is approximable by $\phi({\mr x}\,;{\gr z})$ then for every $n$

\hfil$\displaystyle\A {\gr z_1},\dots,{\gr z_n}\;\E {\mr x}\ \bigwedge^n_{i=1}\big[\phi({\mr x}\,;{\gr z_i})\ \iff\ {\gr z_i}\in{\gr \D}\big]$.


So the same holds for any $\grC\equiv\grD$.
A similar remark apply to approximability from below and from above (e.g.\@ for approximability from below, add the conjunct $\A {\gr z}\,\big[\phi({\mr x}\,;{\gr z})\imp {\gr z}\in\grD\big]$ to the formula above).\QED
\end{remark}

From the following easy observation of Chernikov and Simon~\cite{CS} we obtain an interesting (and misterious) quantifier elimination result originally due to Shelah, see Corollary~\ref{corol_sh_exp_qe} below.

% \begin{proposition}
% Let $\C\subseteq\U^{|{\gr z},w|}$ be approximated from below by the formula $\phi({\mr x}\,;{\gr z},w)$. Then $\grD=\big\{{\gr z}:\E w\ \big({\gr z}\,w\in\C\big)\big\}$ is approximated from below by the formula $\E w\,\phi({\mr x}\,;{\gr z}\,w)$.
% \end{proposition}
% 
% \begin{proof}
% Let $B\subseteq\U$ be finite. 
% We want ${\mr a}\in\U^{|{\mr x}|}$ such that 
% 
% \ceq{\ssf{a.}\hfill \E w\ \big({\gr b}\,w\in\C\big)}{\iff}{\E w\,\phi({\mr a}\,;{\gr b}\,w)}\hfill for every ${\gr b}\in B^{|{\gr z}|}$
% 
% \ceq{\ssf{b.}\hfill \A {\gr z}\;\Big[\E w\,\phi({\mr a}\,;{\gr z}\,w)}{\imp}{\E w\ \big({\gr z}\,w\in\C\big)\Big]}
% 
% Let $C\subseteq\U$ be a finite set such that 
% 
% \ceq{\ssf{c.}\hfill \E w\in C^{|w|}\ \big({\gr b}\,w\in\C\big)}{\iff}{\E w\ \big({\gr b}\,w\in\C\big)}\hfill for every ${\gr b}\in B^{|{\gr z}|}$
% 
% As $\C$ is approximable from below, there is an ${\mr a}$ such that
% 
% \ceq{\ssf{a'.}\hfill {\gr b}\,c\in\C}{\iff}{\phi({\mr a}\,;{\gr b}\,c)}\hfill for every ${\gr b}\,c\in \big(B\cup C\big)^{|{\gr z}\,w|}$
% 
% \ceq{\ssf{b'.}\hfill \A {\gr z}\,w\ \Big[\phi({\mr a}\,;{\gr z}\,w)}{\imp}{{\gr z}\,w\in\C\Big]}
% 
% We obtain \ssf{b} from \ssf{b'} simply by logic. 
% Implication $\imp$ in \ssf{a} follows from \ssf{a'} and \ssf{c}. 
% Implication $\pmi$ follows from \ssf{b}.
% \end{proof}

\begin{proposition}\label{prop_sh_exp_qe}
Let $\C\subseteq\U^{|y\,{\gr z}|}$ be approximated from below by the formula $\phi({\mr x}\,;y\,{\gr z})$.
Then $\grD=\big\{{\gr z}:\E y\ \big(y\,{\gr z}\in\C\big)\big\}$ is approximated from below by the formula $\E y\,\phi({\mr x}\,;y\,{\gr z})$.
\end{proposition}

\begin{proof}
\def\vl{\gr}
Let $B\subseteq\U$ be finite. 
We want ${\mr a}\in\U^{|{\mr x}|}$ such that 

\ceq{\ssf{a.}\hfill \E y\ \big(y\,{\vl b}\in\C\big)}{\iff}{\E y\,\phi({\mr a}\,;y\,{\vl b})}\hfill for every ${\vl b}\in B^{|{\vl z}|}$

\ceq{\ssf{b.}\hfill \A {\vl z}\;\Big[\E y\,\phi({\mr a}\,;y\,{\vl z})}{\imp}{\E y\ \big(y\,{\vl z}\in\C\big)\Big]}

Let $D\subseteq\U$ be a finite set such that 

\ceq{\ssf{c.}\hfill \E y\in D^{|y|}\ \big(y\,{\vl b}\in\C\big)}{\iff}{\E y\ \big(y\,{\vl b}\in\C\big)}\hfill for every ${\vl b}\in B^{|{\vl z}|}$

As $\C$ is approximable from below, there is an ${\mr a}$ such that

\ceq{\ssf{a'.}\hfill d\,{\vl b}\in\C}{\iff}{\phi({\mr a}\,;d\,{\vl b})}\hfill for every $d\,{\vl b}\in \big(D\cup B\big)^{|y\,{\vl z}|}$

\ceq{\ssf{b'.}\hfill \A y\,{\vl z}\ \Big[\phi({\mr a}\,;y\,{\vl z})}{\imp}{y\,{\vl z}\in\C\Big]}

We obtain \ssf{b} from \ssf{b'} simply by logic. 
Implication $\imp$ in \ssf{a} follows from \ssf{a'} and \ssf{c}. 
Implication $\pmi$ follows from \ssf{b}.
\end{proof}



\begin{corollary}
If $p\in S_{{\mr x}}(\U)$ is honestly definable then the family of sets externally definable by $p$ is closed under quantifiers and Boolean combinations. 
\end{corollary}

\begin{proof}
The sets externally definable by $p({\mr x})$ are always closed under Boolean operations. By the proposition above, they are closed under quantifiers. 
\end{proof}



\section{Recycle bin}


\subsection{Ziegler's example}




The language contains a ternary relation which we denote by $x<_zy$ and a unary function $f(x)$. 
\begin{itemize}
  \item[1.] for any fixed $z$, the relation $<_z$ is a dense liner order with least element $z$;
  \item[2.] $x<_zy<_zw\imp y<_xw$;
  \item[3.] $\displaystyle f^nz=z\wedge\bigwedge^{n-1}_{i=0}f^i(z)<_zf^{i+1}(z)$.
\end{itemize}

 
Write $r(x,y)$ for $x=y\vee y<_xf(x)\vee x<_yf(y)$




\end{document}

The language $L$ consists of a binary relation symbol $<$ and, for every positive integer $n$, a binary relation denoted by $|x-y|<n$. Consider the structure with domain $\QQ$ and the relation symbols interpreted in the natural way. Let $T$ be the theory of this structure.

\begin{proposition}
The theory $T$ has elimination of quantifiers.
\end{proposition}




$n\le d(x,z)\ \wedge\ m\le d(x,y)\ \imp\ n+m\le d(y,z)$

$x<y<z\ \wedge\ d(x,y)<n\ \wedge\ d(y,z)<m\ \imp\ d(x,z)<m+n$


$n\le d(x,y) \ \imp\ n+1\le d(x,y)$

$x<y\ \wedge\ n+1\le d(x,y) \ \imp\ \E z\ \big[x<z<y\ \wedge\ n\le d(x,z)\ \wedge\ 1\le d(z,y)\big]$


